\section{Zadanie 3}
Niech $P$ będzie macierzą przejścia skończonego, nieprzywiedlnego i nieokresowego
łańcucha Markowa $(M_t)_{t \ge 0}$.
Załóżmy, że istnieją stałe $c < \tfrac12$ oraz $T \ge 0$ takie, że
\[
\tau(c) \le T.
\]
Wówczas dla każdego $\varepsilon > 0$ zachodzi
\[
\tau(\varepsilon)
\le
\left\lceil \frac{\ln \varepsilon}{\ln(2c)} \right\rceil T.
\]

\subsection*{Rozwiązanie}
Rozważmy dwa dowolne stany początkowe $X_0 = x$ oraz $Y_0 = y$.
Z definicji czasu mieszania $\tau(c)$ otrzymujemy
\[
\|P_x^T - \pi\|_{\mathrm{TV}} \le c
\quad \text{oraz} \quad
\|P_y^T - \pi\|_{\mathrm{TV}} \le c.
\]
Stąd, z nierówności trójkąta,
\[
\|P_x^T - P_y^T\|_{\mathrm{TV}} \le 2c.
\]

Na mocy\\
\medskip
\noindent
\textbf{Lematu o sprzęganiu (Lemat 12.3).}
\emph{
Niech $\sigma_X$ oraz $\sigma_Y$ będą rozkładami prawdopodobieństwa
na skończonej przestrzeni stanów $S$.
Niech $Z=(X,Y)$ będzie dowolną zmienną losową o wartościach w $S\times S$,
taką że $X$ ma rozkład $\sigma_X$, a $Y$ ma rozkład $\sigma_Y$.
Wówczas
\[
\Pr(X \neq Y) \ge \|\sigma_X - \sigma_Y\|_{\mathrm{TV}}.
\]
Ponadto istnieje sprzęgnięcie $Z=(X,Y)$, dla którego zachodzi równość:
\[
\Pr(X \neq Y) = \|\sigma_X - \sigma_Y\|_{\mathrm{TV}}.
\]
}
istnieje sprzężenie
\[
Z_{T,x,y} = (X_T, Y_T),
\]
gdzie $X_T$ ma rozkład $P_x^T$, a $Y_T$ ma rozkład $P_y^T$, takie że
\[
\mathbb{P}(X_T \neq Y_T) \le 2c.
\]

Rozważmy teraz nowy łańcuch Markowa $(M'_t)$ o macierzy przejścia $P^T$.
Odpowiada on wykonywaniu $T$ kroków łańcucha $(M_t)$ naraz.
Dla tego łańcucha, przy powyższym sprzęganiu, jeśli dwie kopie łańcucha
$(M'_t)$ startują ze stanu $(x,y)$, to następny stan pary jest zadany
przez rozkład $Z_{T,x,y}$, co implikuje, że prawdopodobieństwo braku
sprzęgnięcia w jednym kroku wynosi co najwyżej $2c$.

Indukcyjnie otrzymujemy, że prawdopodobieństwo, iż sprzężenie łańcucha
$(M'_t)$ nie nastąpi w ciągu $k$ kroków, jest co najwyżej
\[
(2c)^k.
\]

Z lematu o sprzęganiu wynika zatem, że łańcuch $(M'_t)$ jest w odległości
wariacyjnej co najwyżej $\varepsilon$ od swojego rozkładu stacjonarnego
po $k$ krokach, o ile
\[
(2c)^k \le \varepsilon.
\]
Warunek ten jest spełniony dla
\[
k = \left\lceil \frac{\ln \varepsilon}{\ln(2c)} \right\rceil.
\]

Ponieważ łańcuchy $(M_t)$ oraz $(M'_t)$ mają ten sam rozkład stacjonarny,
a każdy krok łańcucha $(M'_t)$ odpowiada $T$ krokom łańcucha $(M_t)$,
otrzymujemy
\[
\tau(\varepsilon)
\le
\left\lceil \frac{\ln \varepsilon}{\ln(2c)} \right\rceil T.
\]
